\begin{abstract}
% CS197 move 1: problem motivation -- nonresponse and comparator dependence.
Cognitive behavioral therapy (CBT) can benefit patients while leaving substantial nonresponse and uncertainty about its specific mechanisms \citep{Loerinc2015,Carpenter2018}.
% CS197 move 2: assumption -- package benefit does not identify the component.
Evidence for a multicomponent treatment does not identify cognitive restructuring as the necessary source of benefit \citep{Jacobson1996}.
% CS197 move 3: contribution -- broader critical scope requested by the author.
We develop a critical meta-review of reasons CBT can fail across conditions, including anxiety, depression and other clinical applications.
% CS197 move 4: instantiation -- distinguish criticism types and evidence levels.
The protocol separates nonresponse, comparator choice, bias, component mechanisms, executive demands, durability and treatment burden.
Adult attention-deficit/hyperactivity disorder and baseline executive function form one case study.
% CS197 move 5: evaluation -- targeted evidence and incomplete systematic coverage.
The current report provides a cited seed synthesis, a preliminary screen of 21 case-study records, reproducible ADHD search receipts and preserved author drafts.
Cross-condition evidence supports benefits while showing that interpretation depends on the comparator and condition \citep{Cuijpers2025}.
% CS197 move 6: implication -- credible criticism with explicit limits.
A defensible critique must identify the treatment, outcome and population for which a limitation is supported, rather than infer universal ineffectiveness from an uncertain mechanism.
\end{abstract}
